\section{相关工作}
\label{sec:related}

\subsection{检测与跟踪}
\label{sec:related-det-mot}

近年来，3D检测在许多工作中得到了广泛研究，例如基于激光雷达的方法 \cite{yin2021center, chen2022focal, chen2023largekernel3d, chen2023voxelnext, zhang2023fully}、基于相机的方法 \cite{huang2021bevdet, li2022bevformer, liu2022petr, li2023bevdepth, wang2023exploring} 以及多模态融合方法 \cite{bai2022transfusion, liu2023bevfusion, yan2023cross}。由于所提方法主要关注基于相机的方法，因此主要介绍该类方法。由于密集表示的特性，鸟瞰视图（BEV）特征被广泛用于许多优秀工作中。LSS \cite{philion2020lift} 首次提出通过深度预测将图像特征投影到BEV特征上，基于此的一系列后续工作 \cite{huang2021bevdet, li2023bevdepth, li2022bevformer, yang2023bevformer, han2023exploring, li2023bevstereo} 取得了更好的检测性能。另一条3D感知分支采用稀疏路线，无需任何密集表示。DETR3D \cite{wang2022detr3d} 首次提出使用3D查询（Query）聚合多视角2D特征，以稀疏方式实现3D感知。PETR和Sparse4D系列 \cite{liu2022petr, liu2023petrv2, wang2023focal, yan2023cross, zhang2023fully, lin2022sparse4d, lin2023sparse4dv2, lin2023sparse4d} 进一步持续探索基于完全稀疏查询的高效时空聚合方案，在效率和性能上均显示出显著优势。

通常，跟踪从相邻多帧感知结果中生成障碍物轨迹，为下游任务提供输入。传统算法 \cite{weng20203d, yin2021center, wang2023exploring} 采用"基于检测的跟踪（Tracking-by-Detection）"范式，主要关注建立跟踪轨迹与新到来的感知结果之间的关联，其中跟踪性能受检测结果影响显著。受DETR \cite{codevilla2018end} 启发，更多工作 \cite{lin2021global, zeng2022motr, zhang2023motrv2, yu2023motrv3, sun2020transtrack, pang2022simpletrack, meinhardt2022trackformer} 探索了通过在相邻帧间传播查询和监督关联来实现联合检测与跟踪的范式。MUTR3D \cite{zhang2022mutr3d} 将该范式从2D迁移到3D。PF-Track \cite{pang2023standing} 和Sparse4Dv3 \cite{lin2023sparse4d} 进一步强调并有效利用了时空信息，显著提升了性能。然而，在联合训练过程中，不同类型查询之间存在明显的优化不平衡 \cite{yu2023motrv3}，这可能严重损害检测性能。

\subsection{在线建图}
\label{sec:related-mapping}

在自动驾驶中，地图可以提供丰富的场景信息，这对驾驶安全至关重要。基于离线高精地图（HD-Map）的先前解决方案需要大量人力和资源，限制了其可扩展性 \cite{li2022hdmapnet}。近年来，仅使用车载传感器在线构建高精地图引起了广泛关注 \cite{li2022hdmapnet, liu2023vectormapnet, qiao2023end, ding2023pivotnet, yuan2024streammapnet}。HDMapNet \cite{li2022hdmapnet} 通过像素级分割和启发式后处理预测地图。VectorMapNet \cite{liu2023vectormapnet} 提出通过自回归Transformer逐步细化地图元素。BeMapNet \cite{qiao2023end} 以端到端方式构建实例级矢量化高精地图，并采用分段贝塞尔曲线建模地图元素。PivotNet \cite{ding2023pivotnet} 进一步提出基于枢轴（Pivot）表示的地图构建方法。StreamMapNet \cite{yuan2024streammapnet} 采用流式策略，通过时序融合提升性能。然而，上述所有方法均使用密集BEV特征进行环境表示，其计算成本随感知距离呈平方增长，使得在端到端模型中与其他模块联合工作的成本更高。

\subsection{运动预测}
\label{sec:related-forecasting}

在传统自动驾驶系统中，运动预测模块预测周围智能体的未来轨迹或意图，以促进自车的正确决策和安全驾驶。早期方法将驾驶场景渲染为栅格化图像，并通过现成的卷积网络预测未来运动\cite{chai2019multipath, phan2020covernet}。VectorNet \cite{gao2020vectornet} 首次提出用于运动预测的高效矢量化表示，为许多近期方法奠定了基础。基于图的方法也被广泛用于建模和提取智能体与地图之间的相对关系 \cite{zhao2021tnt, varadarajan2022multipath++, zhou2022hivt}。QCNet\cite{zhou2023query}、MTR系列 \cite{shi2022motion, shi2024mtr++} 进一步以更高效的方式探索环境表示并提升预测性能。PnPNet \cite{liang2020pnpnet} 和ViP3D \cite{gu2023vip3d} 在端到端方式下应用轨迹预测方面进行了大量尝试和努力。

\subsection{规划}
\label{sec:related-planning}
作为自动驾驶流程接近末端的关键组件，规划直接影响系统的最终性能。
早期算法采用多种方法来实现规划的特定目标，包括轨迹优化 \cite{diachuk2022motion, shi2022path}、图搜索 \cite{kim2014road, reda2024path}、基于采样的方法 \cite{ma2015efficient, li2020adaptive, orthey2023sampling}。在过去几年中，基于强化学习的方法 \cite{kendall2019learning, aradi2020survey, peng2022drl} 逐渐兴起。然而，由于仿真环境与现实环境之间存在显著差距，这些方法在实际驾驶场景中通常表现不佳。近年来，一些令人印象深刻的工作 \cite{hu2023planning, jiang2023vad} 尝试以端到端方式实现自车的规划。其中，UniAD\cite{hu2023planning} 引入了一个统一的面向规划的框架，在密集BEV中心范式中集成了全栈驾驶任务。VAD \cite{jiang2023vad} 将驾驶场景建模为完全矢量化表示，并利用显式的实例级规划约束来提高规划安全性。
